% ============================================================
% Section 9: Conclusion
% ============================================================

\section{Conclusion}
\label{sec:conclusion}

We surveyed how edge caching models user mobility, classified 24 caching studies by model and by how their decisions use mobility, and explained their conflicting findings through four mechanisms. The question in the title has no single answer. Mobility can help or hurt, and a MATLAB benchmark with seven mobility models shows when.

\emph{Small-cell caching.} Mobility helps a placement that spreads different files over neighboring cells. Its gain over most-popular caching ranges from 0.7\% to 24.3\% across models at the same nominal speed and grows by about 10.4\% for each additional cell a user can reach before the deadline. Mobility hurts reactive caching at the serving cell, which loses exactly the deliveries that a handover interrupts.

\emph{D2D caching.} Mobility helps short transfers and hurts long ones, and in between the effect is non-monotone. The offloading ratio follows from the contact-length distribution through a closed-form expression, with a mean error of 0.004.

\emph{Evaluation choices.} The metric and the user population can reverse a conclusion. The snapshot hit ratio reports a loss where deliveries show a gain, and when most users are static the mobility-aware gain nearly disappears while a placement trained on moving users does worse than most-popular caching.

Results obtained under different models become comparable once a paper reports the reach, the probability of staying in the current cell and the contact-length distribution, together with a delivery-aware hit ratio (Section~\ref{sec:lessons}). The most pressing open problems are trace-calibrated models that can still be swept and bounds on the gain under mobility-prediction error (Section~\ref{sec:open}).
